% Options for packages loaded elsewhere
\PassOptionsToPackage{unicode}{hyperref}
\PassOptionsToPackage{hyphens}{url}
\documentclass[
  ignorenonframetext,
]{beamer}
\newif\ifbibliography
\usepackage{pgfpages}
\setbeamertemplate{caption}[numbered]
\setbeamertemplate{caption label separator}{: }
\setbeamercolor{caption name}{fg=normal text.fg}
\beamertemplatenavigationsymbolsempty
% remove section numbering
\setbeamertemplate{part page}{
  \centering
  \begin{beamercolorbox}[sep=16pt,center]{part title}
    \usebeamerfont{part title}\insertpart\par
  \end{beamercolorbox}
}
\setbeamertemplate{section page}{
  \centering
  \begin{beamercolorbox}[sep=12pt,center]{section title}
    \usebeamerfont{section title}\insertsection\par
  \end{beamercolorbox}
}
\setbeamertemplate{subsection page}{
  \centering
  \begin{beamercolorbox}[sep=8pt,center]{subsection title}
    \usebeamerfont{subsection title}\insertsubsection\par
  \end{beamercolorbox}
}
% Prevent slide breaks in the middle of a paragraph
\widowpenalties 1 10000
\raggedbottom
\AtBeginPart{
  \frame{\partpage}
}
\AtBeginSection{
  \ifbibliography
  \else
    \frame{\sectionpage}
  \fi
}
\AtBeginSubsection{
  \frame{\subsectionpage}
}
\usepackage{iftex}
\ifPDFTeX
  \usepackage[T1]{fontenc}
  \usepackage[utf8]{inputenc}
  \usepackage{textcomp} % provide euro and other symbols
\else % if luatex or xetex
  \usepackage{unicode-math} % this also loads fontspec
  \defaultfontfeatures{Scale=MatchLowercase}
  \defaultfontfeatures[\rmfamily]{Ligatures=TeX,Scale=1}
\fi
\usepackage{lmodern}
\ifPDFTeX\else
  % xetex/luatex font selection
\fi
% Use upquote if available, for straight quotes in verbatim environments
\IfFileExists{upquote.sty}{\usepackage{upquote}}{}
\IfFileExists{microtype.sty}{% use microtype if available
  \usepackage[]{microtype}
  \UseMicrotypeSet[protrusion]{basicmath} % disable protrusion for tt fonts
}{}
\makeatletter
\@ifundefined{KOMAClassName}{% if non-KOMA class
  \IfFileExists{parskip.sty}{%
    \usepackage{parskip}
  }{% else
    \setlength{\parindent}{0pt}
    \setlength{\parskip}{6pt plus 2pt minus 1pt}}
}{% if KOMA class
  \KOMAoptions{parskip=half}}
\makeatother
\setlength{\emergencystretch}{3em} % prevent overfull lines
\providecommand{\tightlist}{%
  \setlength{\itemsep}{0pt}\setlength{\parskip}{0pt}}
% Slide layout defaults for warning-clean scientific decks.
\usepackage{etoolbox}
\IfFileExists{xurl.sty}{\usepackage{xurl}}{}
\IfFileExists{seqsplit.sty}{\usepackage{seqsplit}}{\newcommand{\seqsplit}[1]{#1}}
\protected\def\breakseq#1{\seqsplit{#1}}
\protected\def\breaktt#1{\begingroup\ttfamily\seqsplit{#1}\endgroup}
\setlength{\emergencystretch}{6em}
\tolerance=5000
\hbadness=10000
\hfuzz=1pt
\setlength{\tabcolsep}{2pt}
\AtBeginEnvironment{longtable}{\tiny\renewcommand{\arraystretch}{0.86}\setlength{\tabcolsep}{1pt}}
\AtBeginEnvironment{tabular}{\tiny\renewcommand{\arraystretch}{0.86}\setlength{\tabcolsep}{1pt}}

% Natbib and cross-reference fallbacks — slides don't load natbib
% or cleveref, but combined-PDF manuscript prose may emit these
% commands. The fallback renders citations as a bracketed key list
% and cross-references as detokenized labels so slides don't fail on
% undefined control sequences or raw underscores. \providecommand is
% a no-op if packages load later.
\providecommand{\citep}[1]{[#1]}
\providecommand{\citet}[1]{#1}
\providecommand{\citealp}[1]{#1}
\providecommand{\citeauthor}[1]{#1}
\providecommand{\citeyear}[1]{#1}
\providecommand{\cref}[1]{\texttt{\detokenize{#1}}}
\providecommand{\Cref}[1]{\texttt{\detokenize{#1}}}
\usepackage{bookmark}
\IfFileExists{xurl.sty}{\usepackage{xurl}}{} % add URL line breaks if available
\urlstyle{same}
% fallback for those not using the hyperref driver hyperxmp:
\makeatletter
\@ifundefined{xmpquote}{\newcommand{\xmpquote}[1]{#1}}{}
\makeatother
\hypersetup{
  hidelinks,
  pdfcreator={LaTeX via pandoc}}

\author{\texorpdfstring{}{}}
\date{}

\begin{document}

\section{Introduction}\label{sec:introduction}

\begin{frame}[allowframebreaks]{Bounded Research As Infrastructure}
\protect\phantomsection\label{sec:bounded-research-infrastructure}
AutoResearch systems are most useful when their planning, evidence,
evaluation, and review surfaces remain inspectable. The recent pattern
popularized by bounded coding-agent research loops is simple: define a
tractable objective, try candidate changes under a budget, keep the
result that improves the metric, and leave a replayable trace of what
happened \breakseq{[@karpathy\_autoresearch\_2026]}. That pattern is powerful,
but it is also easy to overstate. A public research template should show
how to run the loop without hiding cost, evidence, review, or execution
boundaries.
\end{frame}

\begin{frame}[fragile,allowframebreaks]{Exemplar Task}
\protect\phantomsection\label{sec:exemplar-task}
This project implements that safer version. The central task is
\breaktt{small\ MNIST\ neural-network\ classification}:
\breaktt{MNIST\ handwritten\ digit\ database} from the handwritten-digit
database \breakseq{[@lecun\_mnist\_database]}, a \breaktt{nearest-centroid}
baseline, and a finite list of candidate model families
(\breaktt{MLP,\ nearest-centroid,\ softmax\ regression,\ tiny\ patch-attention}).
The AutoResearch loop is responsible for proposing candidate
configurations, evaluating them against \texttt{test\_accuracy},
selecting the best result with deterministic tie-breaking, and writing
the evidence needed to review the claim.
\end{frame}

\begin{frame}[fragile,allowframebreaks]{Contribution And Boundary}
\protect\phantomsection\label{sec:contribution-boundary}
The contribution is not a new \texttt{MNIST} classifier. It is a
template-level demonstration of how bounded AutoResearch can be
orchestrated through the same lifecycle used for reproducible papers:
tests run first, analysis writes structured artifacts, rendering
hydrates manuscript variables, validation checks evidence and readiness,
and copy stages publish final deliverables. The default path makes no
network calls, no LLM calls, executes no generated code, and never
treats a generated review packet as human approval. The methods in
@sec:methodology define that boundary, and the results in @sec:results
report only what the validated artifacts support.
\end{frame}

\begin{frame}[allowframebreaks]{Process Organization Motif}
\protect\phantomsection\label{sec:process-organization-motif}
The exemplar also treats research orchestration itself as a first-class
research object. In that limited operational sense, ``research for
research's sake'' means that programs, ledgers, evidence registries,
validation gates, and manuscript hydration do more than report a result:
they maintain the conditions under which the next research step can be
inspected, replayed, criticized, and extended. This is a process
analogy, not a moral claim that software is an end in itself or a
biological claim that the template is alive.
\end{frame}

\begin{frame}[fragile,allowframebreaks]{Related Work And Current Trends}
\protect\phantomsection\label{sec:related-work}
\begin{block}{Information Overload And Machine-Readable Science}
\protect\phantomsection\label{sec:information-overload-machine-readable-science}
The current AutoResearch literature is driven by a practical bottleneck:
scientific output is growing faster than document-centered review and
synthesis can absorb. Recent science-of-science evidence complicates any
simple productivity story: AI-augmented scientists can publish and be
cited more often, but the same adoption pattern can narrow the
collective range of topics and interactions in science
\breakseq{[@hao\_ai\_tools\_focus\_2026]}. This is a 2026 Nature result, not a
2024 analysis, and it motivates governance rather than celebration.

One response is to make literature synthesis more source governed.
OpenScholar uses retrieval-augmented generation over a large scientific
passage store and reports citation accuracy improvements on literature
synthesis tasks \breakseq{[@asai\_openscholar\_2026]}. PaperQA2 similarly
evaluates literature-search agents against expert scientific tasks and
emphasizes cited answers, contradiction detection, and factuality
\breakseq{[@skarlinski\_paperqa2\_2024]}. STORM, PaperQA, and GPT Researcher
are adjacent source-grounded writing systems that motivate this
project's insistence on citation-backed claims and visible evidence
surfaces \breakseq{[@shao\_storm\_2024; @lala\_paperqa\_2023;
@gpt\_researcher\_2026]}. The common lesson is that automated writing
is not enough: claims must remain tied to inspectable sources,
artifacts, and evaluation records.
\end{block}

\begin{block}{Structured Distillation And Conceptual Knowledge
Substrates}
\protect\phantomsection\label{sec:structured-distillation-knowledge-substrates}
The Discovery Engine proposes a more structural answer to the same
overload problem. It distills publications into source-linked knowledge
artifacts, organizes those artifacts under a conceptual schema, encodes
them into a high-dimensional Conceptual Tensor, and unrolls that tensor
into graph and vector views for agent navigation
\breakseq{[@baulin\_discovery\_engine\_2025]}. This project does not construct
a Conceptual Nexus Model or claim domain-scale literature synthesis. It
adopts a much smaller analogue: outputs, ledgers, evidence registries,
figure records, and hydrated manuscript variables are file-backed
objects whose provenance can be checked before publication.

The representational background is broader than one framework. FAIR
principles argue for data that are findable, accessible, interoperable,
and reusable by people and machines \breakseq{[@wilkinson2016fair]}. RO-Crate
and Workflow Run RO-Crate package research artifacts and computational
executions with linked-data metadata \breakseq{[@soiland\_reyes2022rocrate;
@soiland\_reyes2024workflow\_run\_rocrate]}. Hyperdimensional computing
and vector symbolic architectures provide one route for robust
high-dimensional symbolic-subsymbolic representations
\breakseq{[@heddes\_hdc\_2024]}, while tensor factorization methods such as
TuckER and mixed-geometry tensor factorization show how multi-relational
knowledge graphs can be completed and queried as structured tensors
\breakseq{[@balazevic\_tucker\_2019; @yusupov\_mixed\_geometry\_2025]}. The
local contribution here is not a new knowledge representation method; it
is an executable template that makes a small research workflow
compatible with that machine-readable direction.
\end{block}

\begin{block}{End-To-End AutoResearch Systems}
\protect\phantomsection\label{sec:end-to-end-autoresearch-systems}
The most ambitious AutoResearch systems now aim at the entire scientific
lifecycle. The AI Scientist assembles idea generation, experiment
execution, paper writing, and automated review
\breakseq{[@lu\_aiscientist\_2024]}, and its Nature version reports an
end-to-end AI research pipeline whose generated manuscript passed a
workshop peer-review round \breakseq{[@lu\_ai\_scientist\_nature\_2026]}. AI
Scientist-v2 removes more hand-authored scaffolding and uses agentic
tree search for broader hypothesis exploration
\breakseq{[@yamada\_ai\_scientist\_v2\_2025]}. FutureHouse's platform exposes
specialized scientific agents for literature search, deep review,
novelty checking, and chemistry planning
\breakseq{[@futurehouse\_platform\_2025]}, while Robin integrates literature
and data-analysis agents in a lab-in-the-loop discovery workflow
\breakseq{[@ghareeb\_robin\_2026]}.

Survey work is already separating reliable assistance from risky
autonomy. The AI for Auto-Research roadmap describes the full lifecycle
from creation to dissemination, but stresses that novelty,
research-level implementation, and judgment remain fragile under
automation \breakseq{[@kong\_ai\_auto\_research\_2026]}. EXHYTE frames
discovery as an iterative Exploration, Hypothesis generation, and
Testing loop, clarifying where current systems are mature and where
closed-loop autonomy remains thin \breakseq{[@hasib\_exhyte\_2025]}. This
exemplar therefore takes the opposite stance from full autonomy: it
implements a bounded local loop whose candidate space, data, cost,
outputs, and review gates can be audited.
\end{block}

\begin{block}{Autoformalization And Verification}
\protect\phantomsection\label{sec:autoformalization-verification}
Autoformalization supplies a different kind of boundary: instead of only
asking whether generated text is plausible, it asks whether an informal
statement can be translated into a form that a proof assistant or
compiler can check. AlphaProof shows the power of reinforcement learning
over formal mathematical proof search \breakseq{[@hubert\_alphaproof\_2025]}.
Process-driven autoformalization in Lean uses compiler feedback as a
precise signal for improving translations from natural-language
mathematics to formal statements and proofs
\breakseq{[@lu\_process\_driven\_autoformalization\_2024]}. APOLLO turns Lean
feedback into an iterative proof-repair workflow in which generated
proofs are decomposed, patched, and reverified
\breakseq{[@apollo\_lean\_collaboration\_2025]}.

This project does not perform theorem proving, proof repair, or formal
mathematical verification. It borrows the architectural lesson:
generated research artifacts should be checked by deterministic tools
with explicit error surfaces. Here those tools are test suites, schema
checks, evidence registries, render validation, source hygiene greps,
and review packets rather than Lean, Isabelle, or Coq.
\end{block}

\begin{block}{ML For ML And Evolutionary Algorithm Discovery}
\protect\phantomsection\label{sec:ml-for-ml-evolutionary-discovery}
ML-for-ML systems optimize models, code, or algorithms with search loops
that are themselves subject to evaluation. Karpathy's
\texttt{autoresearch} repository frames a minimal version of this
pattern as a prompt-controlled system with a fixed budget, editable code
surface, and comparable metric \breakseq{[@karpathy\_autoresearch\_2026]}.
MLAgentBench and MLE-bench package machine learning tasks as scored,
replayable environments with logs and grading outputs
\breakseq{[@huang\_mlagentbench\_2023; @chan\_mle\_bench\_2024]}.

At a larger scale, AlphaEvolve couples language-model proposals to
evolutionary program search and automated evaluators, producing
algorithmic improvements in mathematics and computing
\breakseq{[@romera\_paredes\_alphaevolve\_2025]}. DeepEvolve adds external
retrieval, multi-file code editing, and debugging to the same basic
proposal-implementation-evaluation loop \breakseq{[@deepevolve\_2025]}. This
manuscript's candidate search is intentionally smaller: a finite list of
configured model families is evaluated against \texttt{test\_accuracy},
with deterministic selection and recorded deferrals rather than
unbounded code mutation.
\end{block}

\begin{block}{Agentic Science, Graph Retrieval, And Epistemic Foraging}
\protect\phantomsection\label{sec:agentic-science-graph-retrieval}
Agentic science surveys describe systems that move from tool use toward
scientific agency across perception, knowledge representation, planning,
experimentation, analysis, and communication
\breakseq{[@wei\_agentic\_science\_2025; @gridach\_agentic\_science\_2025]}.
GraphRAG work adds structured retrieval to that picture: graph
construction and graph-aware retrieval can support multi-hop reasoning,
but benchmarks also show that knowledge-graph RAG remains brittle when
relevant knowledge is incomplete \breakseq{[@graphrag\_bench\_2026;
@brink\_kg\_rag\_2026]}. Active-inference perspectives make a similar
design demand in different language: scientific agents need persistent
uncertainty-aware memory, causal models, counterfactual exploration,
deterministic validation, and human judgment as an architectural
component \breakseq{[@active\_inference\_science\_2025]}.

This exemplar uses those ideas as constraints, not as capabilities it
already possesses. It does not run live literature mining, autonomous
proof search, external agent swarms, graph-based hypothesis hunting, or
self-approval. Instead it exposes bounded candidates, explicit budgets,
local evidence links, local \texttt{MNIST} input data, benchmark-style
scoring, source-linked figures, manuscript hydration, and human review
gates. That is the intended safe baseline for a public template.
\end{block}

\begin{block}{Benchmark, Documentation, And Process Analogies}
\protect\phantomsection\label{sec:benchmark-documentation-process-analogies}
\texttt{MNIST} and LeNet remain useful here because they provide a
compact historical benchmark for small neural networks and handwriting
recognition \breakseq{[@lecun\_gradient\_1998; @lecun\_mnist\_database]}.
Vision Transformers introduce the patch-token pattern for image
classification at scale \breakseq{[@dosovitskiy\_vit\_2020]}; this exemplar
borrows only the patching and attention representation through
\breaktt{Tiny\ patch-attention\ classifier}, then keeps the
implementation inside the configured candidate budget. MLPerf Tiny and
OpenML motivate explicit task descriptions, fixed inputs,
machine-readable run metadata, and checkable metrics
\breakseq{[@banbury\_mlperf\_tiny\_2021; @vanschoren\_openml\_2014]}.
Machine-learning reproducibility checklists motivate reporting data,
seeds, model sizes, hyperparameters, and compute boundaries
\breakseq{[@pineau\_reproducibility\_2020]}.

Dataset and model documentation work further informs the safe boundary.
Datasheets for Datasets motivate explicit reporting of dataset
motivation, composition, collection, and recommended use
\breakseq{[@gebru2021datasheets]}. Model Cards motivate structured reporting of
model context, intended use, evaluation procedure, and limitations
\breakseq{[@mitchell2019model\_cards]}. The diagnostic layer follows the same
conservative reporting posture: calibration is treated as separate from
accuracy \breakseq{[@guo2017calibration]}, binomial accuracy intervals use
Wilson-style score intervals \breakseq{[@wilson1927probable\_inference]},
matched classifier comparison is summarized through paired discordance
\breakseq{[@dietterich1998statistical\_tests]}, deterministic bootstrap
intervals are local resampling diagnostics \breakseq{[@efron1993bootstrap]},
and probability quality is reported with Brier score, negative log
likelihood, and chance-corrected agreement \breakseq{[@brier1950verification;
@cohen1960coefficient]}.

The process language is borrowed cautiously from teleology and
theoretical biology. Kant's account of organized beings treats a natural
purpose as a whole whose parts and whole mutually condition one another
\breakseq{[@ginsborg\_kant\_aesthetics\_teleology]}. Autopoiesis characterizes
living systems through self-producing organization
\breakseq{[@varela\_autopoiesis\_1974]}, and later work connects Kantian
natural purpose to autopoietic individuality
\breakseq{[@weber\_varela\_life\_after\_kant\_2002]}. Moreno and Mossio's
account of biological autonomy emphasizes organizational closure and
self-maintenance \breakseq{[@moreno\_mossio\_biological\_autonomy\_2015]}. This
paper uses those ideas only as disciplined analogies for configured
scientific workflows whose artifacts help reproduce, constrain, and
evaluate subsequent artifacts.

Security and supply-chain references enter with the same restraint.
NIST's zero-trust architecture treats verification as explicit and
continuous rather than inherited from a trusted perimeter
\breakseq{[@nist\_sp800\_207\_zero\_trust]}. The NIST Secure Software
Development Framework emphasizes repeatable practices for reducing
software vulnerability risk \breakseq{[@nist\_sp800\_218\_ssdf]}. SLSA frames
software-artifact provenance and supply-chain integrity as a graded
assurance problem \breakseq{[@slsa\_spec\_latest]}, while MITRE ATT\&CK T1195
names supply-chain compromise as a concrete adversary technique
\breakseq{[@mitre\_attack\_t1195]}. This exemplar borrows those frameworks as
disciplined analogies for local research artifact integrity: checksums,
inventories, review gates, and explicit non-claims, not production
deployment certification.
\end{block}
\end{frame}

\end{document}
